\chapter{Estado da Arte}\label{cap:estArte}

A corrida autônoma ocupa uma posição peculiar na literatura de veículos
autônomos: é o único cenário em que o objetivo declarado é operar o veículo no
limite da aderência, e em que o critério de sucesso --- o tempo de volta --- é
um número escalar, medido sem ambiguidade e comparável entre trabalhos
\cite{Betz2022Survey}. Essa combinação transformou a corrida num banco de
provas para percepção, planejamento, controle e aprendizado de máquina, num
regime em que as aproximações que sustentam a condução urbana deixam de valer.

Este capítulo organiza a literatura em torno das decisões de projeto que este
trabalho de fato tomou. A cascata implementada tem três etapas encadeadas, e
cada uma delas herda um problema aberto de uma linha de pesquisa distinta:

\begin{itemize}
    \item \textbf{E1} --- um especialista de \textit{pure pursuit} operando
          sobre uma linha de corrida de curvatura mínima, que fixa o tempo de
          referência $T_{pp}$ e produz os rótulos das etapas seguintes. Herda a
          literatura de otimização de trajetória e seguimento geométrico.
    \item \textbf{E2} --- uma rede convolucional unidimensional compacta que
          aprende por imitação a política do especialista, com ruído de coleta
          estilo DART calibrado por pista. Herda a literatura de clonagem de
          comportamento e o problema de desvio de distribuição.
    \item \textbf{E3} --- um agente de aprendizado por reforço \textit{residual}
          cuja ação se soma à saída da rede de imitação, e não à do especialista
          geométrico. Herda a literatura de aprendizado por reforço para corrida
          e de políticas residuais.
\end{itemize}

Uma advertência precede tudo o que se segue. \textbf{Todos os resultados deste
trabalho citados neste capítulo foram obtidos em simulação}, no ambiente
\texttt{f110\_gym} \cite{OKelly2020F1TENTH}. Sempre que este capítulo comparar um
resultado próprio com um resultado da literatura obtido em hardware, a
assimetria será dita explicitamente.

Além dessas três etapas, o comparativo inclui outros dois tipos de piloto, que
completam o conjunto. O primeiro são dois controladores reativos clássicos, o
Disparity Extender e o seguidor de gap por descontinuidade, que decidem o comando
só a partir da varredura do LiDAR, sem pose, sem mapa e sem linha de corrida; os
dois são descritos na Seção~\ref{sec:arte-classico}, entre os métodos reativos.
Ao contrário deles, o \textit{pure pursuit} do E1 usa a pose fornecida pelo
simulador e a linha de corrida calculada sobre o mapa. O segundo é
o piloto humano, que conduz o veículo simulado com um \textit{joystick} sob os
mesmos limites de ação das políticas e em duas condições de informação: a pista
inteira vista de cima e apenas os pontos do LiDAR. A Seção~\ref{sec:arte-humanos}
reúne o que a literatura diz sobre essa referência humana. Assim, a lista E1--E3
acima enumera as etapas do pipeline de aprendizado, e não todos os pilotos
medidos.

O capítulo está organizado em oito partes. A Seção~\ref{sec:arte-plataformas}
trata das plataformas e simuladores para corrida em escala reduzida, e do que
significa escolher um simulador em vez de outro. A
Seção~\ref{sec:arte-classico} percorre o planejamento e o controle clássicos:
linha de corrida, \textit{pure pursuit}, controle preditivo e os métodos
reativos. A Seção~\ref{sec:arte-il} trata do aprendizado por imitação e do
desvio de distribuição que o DART ataca. A Seção~\ref{sec:arte-rl} trata do
aprendizado por reforço aplicado à corrida, com atenção particular à
formulação residual. A Seção~\ref{sec:arte-rand} trata, em um parágrafo, da
randomização de domínio, que justifica a variação de condições usada no
treino e na avaliação. A Seção~\ref{sec:arte-timetrial} descreve o formato
\textit{time trial}, que o comparativo imita, a
Seção~\ref{sec:arte-humanos} trata de pilotos humanos e controladores
autônomos em corrida, e a Seção~\ref{sec:arte-lacuna}
fecha posicionando a lacuna: o que este trabalho faz que os citados não fazem,
e o que eles fazem que este não faz.

\section{Plataformas e simuladores para corrida autônoma em escala}\label{sec:arte-plataformas}

\subsection{Do veículo de competição ao modelo em escala}

A pesquisa em corrida autônoma se distribui por uma faixa larga de escalas. No
extremo superior estão os veículos de competição de tamanho real, como os da
Indy Autonomous Challenge, cujo software integra percepção multissensorial,
planejamento em ambiente adversarial e controle a mais de 250~km/h
\cite{Betz2023TUM}. No extremo inferior estão os modelos 1:43 usados para
estudar controle ótimo em ciclo fechado, em que a pista cabe numa mesa e a
captura de movimento externa fornece pose exata a alta frequência
\cite{Liniger2015Optimization}. Entre os dois, a escala 1:10 se estabeleceu como
o compromisso dominante: reproduz a dinâmica de um veículo Ackermann com
deslizamento lateral relevante, admite os mesmos sensores e a mesma pilha de
software de um veículo real, e cabe num orçamento acadêmico
\cite{Betz2022Survey}.

A revisão sistemática de Betz e colaboradores \cite{Betz2022Survey}, que examina
233 trabalhos publicados até 2021, é a referência de mapa do campo utilizada
neste capítulo. Ela organiza a literatura em percepção, planejamento, controle e
abordagens fim a fim, e registra dois achados que interessam diretamente ao
posicionamento deste trabalho: a predominância de estudos de veículo único sobre
estudos multiagente, e a carência de conjuntos de dados públicos para o regime de
alta velocidade --- carência que empurra a validação inicial para a simulação e
torna a fidelidade do simulador uma questão de primeira ordem, e não de
engenharia acessória.

\subsection{A plataforma F1TENTH}

A plataforma F1TENTH foi apresentada como ambiente aberto de avaliação para
controle contínuo e aprendizado por reforço \cite{OKelly2020F1TENTH}. Sua
proposta original combina três elementos que raramente aparecem juntos: um
veículo físico de custo acessível com sensoriamento LiDAR e computação
embarcada; um simulador com modelo de veículo de faixa única (\textit{single
track}) calibrado contra o veículo físico; e um conjunto de \textit{benchmarks}
padronizados --- planejamento cinodinâmico estático, planejamento cinodinâmico
\textit{online} e jogos dinâmicos multiagente.

O número mais citado desse trabalho é a discrepância de 0,61\% entre tempos de
volta simulados e medidos no veículo real \cite{OKelly2020F1TENTH}. Ele merece
leitura cuidadosa, porque é frequentemente lido como se o simulador fosse fiel
por construção. Ele não é: a discrepância foi obtida com identificação
sistemática de parâmetros e com randomização de domínio aplicada exatamente para
fechar essa diferença. O que o número atesta é que a diferença \textbf{pode} ser
fechada nessa escala, com esse esforço --- não que ela seja pequena por padrão.

A plataforma consolidou-se também como veículo pedagógico. Agnihotri e
colaboradores \cite{Agnihotri2020Teaching} descrevem o currículo de oito
laboratórios construído sobre o carro 1/10, e Ueter e colaboradores
\cite{Ueter2020CPS} relatam o caso de um projeto estudantil de dois semestres em
sistemas ciber-físicos organizado sobre a mesma base. Esses dois trabalhos são
relevantes aqui menos pelos seus resultados e mais por um efeito colateral: o uso
educacional consolidou um repertório de algoritmos de referência --- seguimento
de parede, seguimento de vão, \textit{pure pursuit} --- que qualquer trabalho
sobre a plataforma pode assumir como conhecido, e contra o qual qualquer
proposta nova é implicitamente comparada.

\subsection{Simuladores: o compromisso entre fidelidade e taxa de amostragem}

Existem pelo menos duas famílias de simuladores para a plataforma, e a escolha
entre elas não é neutra.

A primeira é a dos simuladores baseados em ROS e Gazebo, exemplificada pelo
\textit{f1tenth.dev} \cite{Babu2020Simulator}. A vantagem é a fidelidade da
interface: os mesmos nós, os mesmos tópicos e as mesmas mensagens do veículo
real, o que torna a transferência de código para o hardware quase trivial. A
desvantagem é o custo por passo de simulação --- física tridimensional com
detecção de colisão de malha e um grafo de computação distribuído entre
processos. Para aprendizado por reforço, que exige centenas de milhares a
milhões de transições, esse custo é proibitivo.

A segunda é a dos simuladores dedicados de duas dimensões, com modelo de faixa
única e varredura LiDAR sintetizada por lançamento de raios sobre um mapa de
ocupação. É o caso do ambiente \texttt{f110\_gym} distribuído com a plataforma
\cite{OKelly2020F1TENTH}, adotado por este trabalho. O ganho é de ordem de
grandeza em vazão; a perda é a interface, que deixa de ser ROS2.

A escolha deste trabalho é o ambiente rápido, que permite treinar e avaliar
com o orçamento de amostras que o aprendizado por reforço exige. A interface
ROS2 do veículo não faz parte do escopo: o trabalho compara pilotos dentro do
mesmo simulador e sob o mesmo avaliador, e a fidelidade de interface que a
primeira família oferece não é necessária a essa comparação.

\subsection{Verificação e teste como parte da plataforma}

Duas linhas de trabalho tratam a plataforma não como ambiente de execução, mas
como objeto de verificação.

Ivanov e colaboradores \cite{Ivanov2020Verifying} tomam um controlador neural
treinado por aprendizado por reforço profundo (DDPG e TD3) sobre varreduras
LiDAR e tentam verificar formalmente propriedades de segurança com a ferramenta
Verisig. O resultado é um limite: as ferramentas da época escalam até cerca de
40 raios LiDAR, e a verificação exige um modelo de observação analítico do
sensor. O trabalho também documenta padrões de falha do LiDAR real que podem ser
reproduzidos em simulação, e estabelece um \textit{benchmark} deliberadamente
difícil para a comunidade de verificação. Para este trabalho, a leitura relevante
é negativa e deve ser dita: a rede de imitação aqui empregada consome
\textbf{1080} feixes e é uma rede convolucional, duas características que a
colocam fora do alcance das ferramentas descritas naquele estudo.

Bak e colaboradores \cite{Bak2021Stress} atacam o mesmo problema por amostragem
em vez de por prova. O método usa RRT \cite{Karaman2011Optimal} para explorar
sistematicamente perturbações no cenário de ultrapassagem, e encontra até 4,7
vezes mais colisões que a busca aleatória, chegando a um aumento de 109,6 vezes
na segunda metade da pista. É um argumento forte a favor de teste dirigido
contra teste por amostragem uniforme, e um recado sobre o significado de uma
taxa de conclusão de voltas medida em condições nominais: ela mede o
comportamento médio, não a existência de um cenário adversarial que o quebre.
Este trabalho mede taxa de conclusão de voltas em condições nominais e sob
randomização de coeficiente de atrito, e \textbf{não} executa busca adversarial
--- uma limitação declarada na Seção~\ref{sec:arte-lacuna}.

\section{Planejamento e controle clássicos em corrida}\label{sec:arte-classico}

\subsection{Linha de corrida: da centerline à curvatura mínima}

O problema de encontrar a trajetória mais rápida dentro dos limites da pista é
distinto do problema de planejamento de movimento em ambiente com obstáculos. A
análise formal de Karaman e Frazzoli \cite{Karaman2011Optimal} estabeleceu que
PRM e RRT, embora probabilisticamente completos, \textbf{não} são
assintoticamente ótimos, e propôs PRM*, RRG e RRT* com garantia de otimalidade
assintótica. Esse arcabouço é a referência para planejamento em espaços de
configuração de alta dimensão e ambientes desestruturados, e é justamente por
isso que ele não é a ferramenta natural para a linha de corrida: o circuito é
fechado, periódico e sem obstáculos, e o que se quer é um ótimo global de um
funcional de tempo sobre uma faixa, não um caminho viável num espaço cheio.

A formulação que domina a literatura de corrida é a otimização de curvatura
mínima de Heilmeier e colaboradores \cite{Heilmeier2020MinCurv}, que resolve um
problema quadrático sobre o afastamento lateral em relação à linha central,
minimizando a curvatura ao longo do circuito, e em seguida computa um perfil de
velocidade compatível com um modelo de aderência. A implementação de referência
desse método \cite{TUMFTMRaceTrajectory} é a origem das linhas de corrida
distribuídas com os mapas oficiais da plataforma \cite{F1TENTHRaceTracks}.
Abordagens alternativas usam otimização sem derivadas: o relatório de Misra,
Zheng e Tang \cite{Misra2024Report} emprega CMA-ES para otimizar conjuntamente
a trajetória e o perfil de velocidade sobre a mesma plataforma.

Este trabalho \textbf{não} reimplementa o otimizador. Ele consome a linha de
corrida publicada com os mapas e aplica sobre ela duas correções que a
literatura de otimização de trajetória não precisa fazer, e que aparecem quando
a linha encontra um simulador de dinâmica diferente do que a produziu. A
primeira é geométrica: a linha publicada usa a pista inteira, aproximando-se a
0,925~m do centro onde a meia-largura é 1,10~m; descontada a meia-largura do
veículo, restam cerca de dois centímetros de folga. A linha foi otimizada para
um veículo pontual, e segui-la com perfeição já raspa o muro. A segunda é
dinâmica: o perfil de velocidade foi otimizado para o modelo de pneu do
otimizador, não para o modelo de faixa única deste simulador, e chega a pedir
8~m/s numa curva de raio 6,5~m. As duas correções --- encolher o afastamento
lateral e limitar a velocidade por uma aceleração lateral admissível --- são
calibrações que o modelo mínimo não enxerga. O registro dessa incompatibilidade
é, em si, um dado sobre a portabilidade de linhas de corrida publicadas entre
simuladores.

\subsection{Seguimento geométrico de trajetória: pure pursuit}

O \textit{pure pursuit} \cite{Coulter1992PurePursuit} é o controlador de
seguimento de caminho mais antigo e mais usado da robótica móvel de tração
Ackermann. A formulação é puramente geométrica: escolhe-se um ponto da
trajetória a uma distância de antecipação à frente do veículo e computa-se o
arco de circunferência que liga a pose atual a esse ponto, do qual sai
diretamente o ângulo de esterço. A distância de antecipação é tipicamente afim
na velocidade, o que dá ao controlador um comportamento estável em regime e uma
tendência conhecida a cortar curvas fechadas quando a antecipação é longa
demais.

As limitações do método são bem documentadas e relevantes aqui. Ele não tem
modelo dinâmico, portanto não antecipa perda de aderência; ele responde à
geometria e não ao erro de curvatura, portanto acumula desvio em curvas de raio
variável; e, sobretudo, ele exige \textbf{pose global} --- o veículo precisa
saber onde está no mapa para escolher o ponto de mira. Em um veículo real isso
significa depender de uma pilha de localização, com todo o custo computacional e
todos os modos de falha associados \cite{fox2003adapting, thrun2005}.

Este trabalho usa o \textit{pure pursuit} como \textbf{especialista}, e não como
controlador de entrega. A escolha é motivada por três propriedades: o
controlador é determinístico, é barato e é sem estado interno --- o que, como o
Capítulo~\ref{cap:metodologia} detalha, torna-o imune a uma classe de erro de
avaliação que atingiu as políticas aprendidas. Sobre a linha de corrida
encolhida e com o limite de aceleração lateral, ele completa 50 de 50 voltas sem
colisão e fixa o tempo de referência $T_{pp} = 59{,}000$~s, que é o denominador
de todos os critérios das etapas seguintes.

\subsection{Controle preditivo e suas variantes com aprendizado}

A família de controle preditivo baseado em modelo (MPC) é a resposta clássica ao
que falta ao \textit{pure pursuit}: um modelo do veículo e um horizonte de
predição sobre o qual otimizar. Verschueren e colaboradores
\cite{verschueren2016time} formulam a corrida como problema de tempo mínimo
resolvido por MPC não linear com Hessiana exata; Liniger, Domahidi e Morari
\cite{Liniger2015Optimization} aplicam MPC de contorno a veículos 1:43 em
competição autônoma. O custo dessa família é sempre o mesmo: exige um modelo
identificado, um estimador de estado e um solucionador que feche dentro do
período de controle.

Três linhas atacam esse custo por caminhos diferentes, e todas são pertinentes
ao desenho deste trabalho.

A primeira é o MPC iterativo de aprendizado (LMPC), em que o conjunto seguro e a
função valor terminal são construídos a partir das voltas já executadas, de modo
que o desempenho melhora a cada iteração \cite{Rosolia2020LMPC}. Wang, Zeng e
Mangharam \cite{Wang2020LMPC} aplicam essa formulação à plataforma F1TENTH,
implementando o controlador em C++ sobre o solucionador OSQP e reportando
redução do tempo de volta ao longo das iterações; a validação foi conduzida
primariamente em simulação, e os próprios autores registram a ausência de
validação no veículo real como limitação.

A segunda é a correção do modelo por aprendizado. Kabzan e colaboradores
\cite{Kabzan2019LearningMPC} usam processos gaussianos para corrigir o modelo
nominal do veículo dentro do laço de MPC. Jain e colaboradores, no BayesRace
\cite{Jain2020Bayesrace}, levam a ideia a um extremo instrutivo: partem de um
modelo cinemático estendido com apenas três parâmetros facilmente mensuráveis,
coletam dados com um controlador simples de \textit{pure pursuit}, treinam
processos gaussianos para corrigir a discrepância entre esse modelo e a dinâmica
real, e alimentam com o resultado um MPC não linear. O desempenho se aproxima do
de controladores que usam modelos dinâmicos completos e precisamente
identificados, com esforço de identificação drasticamente menor. A limitação
declarada é o tempo de execução do controlador, da ordem de 0,26~s por passo.

O BayesRace merece destaque especial neste capítulo porque ele é o vizinho mais
próximo da ideia central deste trabalho, e a diferença entre os dois é exatamente
o que dá conteúdo à proposta. Nos dois casos há uma \textbf{correção aprendida
somada a uma base projetada à mão}. No BayesRace o resíduo vive no
\textbf{espaço do modelo}: o processo gaussiano corrige a predição de dinâmica, e
quem decide a ação continua sendo o otimizador. Neste trabalho o resíduo vive no
\textbf{espaço da ação}: não há modelo a corrigir nem otimizador no laço, e a
correção incide diretamente sobre o comando enviado ao atuador. A consequência
prática é de custo computacional e de dependência: o BayesRace precisa de estado
estimado e de solucionador em tempo real; a política composta aqui é uma
avaliação direta de duas redes sobre a varredura LiDAR.

A terceira linha é a substituição do otimizador por uma política aprendida, e é
o assunto das duas seções seguintes.

\subsection{Métodos reativos: seguimento de parede e de vão}

No extremo oposto do MPC estão os controladores reativos, que mapeiam a
varredura do sensor diretamente em comando sem construir mapa, estimar pose ou
planejar trajetória. O seguimento de parede regula a distância lateral a uma
parede de referência por realimentação proporcional-integral-derivativa; o
seguimento de vão, formalizado por Sezer e Gokasan
\cite{Sezer2012FollowGap}, identifica a maior abertura livre no campo de visão e
dirige para o seu centro, ponderando o alvo pela direção do objetivo. O
Disparity Extender, divulgado por Otterness \cite{Otterness2019Disparity} a
partir do carro da equipe da University of North Carolina que venceu a
competição F1TENTH de 2019, é uma variante do seguimento de vão: antes de
escolher a abertura, procura as descontinuidades (disparidades) entre feixes
vizinhos e estende o alcance mais curto de cada uma pela largura do carro, de
modo que o alvo escolhido não passe rente a uma borda; o carro segue então para o
ponto livre mais distante, com velocidade dada pelo espaço à frente. Não há, até
onde foi possível verificar, publicação formal do método; a fonte é a descrição
do autor.

Os três são parte do repertório da plataforma
\cite{Agnihotri2020Teaching, Babu2020Simulator}. Neste trabalho, o seguimento de
vão e o Disparity Extender estão implementados no simulador, conforme descrito
no Capítulo~\ref{cap:proposta}. Eles não competem em
tempo de volta com uma linha de corrida otimizada --- não têm informação para
isso, já que enxergam apenas o horizonte do sensor. O que eles têm, e que
nenhum método baseado em trajetória global tem, é a independência de pose: um
seguidor de vão não depende de mapa nem de localização para operar, e por isso
pode ser levado a um circuito desconhecido.

Essa propriedade é o fio que liga esta seção às seguintes. A pergunta que
organiza a segunda metade do capítulo é se é possível obter o desempenho de uma
trajetória globalmente otimizada mantendo a independência de pose de um
controlador reativo. O aprendizado por imitação é a resposta que a literatura dá
a essa pergunta, e é a etapa E2 deste trabalho.

\section{Aprendizado por imitação e o desvio de distribuição}\label{sec:arte-il}

\subsection{Clonagem de comportamento fim a fim}

A ideia de treinar uma rede para reproduzir os comandos de um condutor a partir
da leitura bruta de um sensor é antiga: o ALVINN \cite{Pomerleau1989ALVINN}
dirigia um veículo a partir de imagens de baixa resolução com uma rede de uma
camada oculta. A retomada moderna \cite{Bojarski2016EndToEnd} demonstrou que uma
rede convolucional treinada com pares imagem-esterço coletados de um condutor
humano aprende a manter faixa sem que nenhum módulo de percepção intermediário
seja especificado.

Transposta para a corrida em escala, a formulação troca a câmera pelo LiDAR
bidimensional, e o condutor humano por um controlador de referência. O
especialista pode ser um \textit{pure pursuit} sobre uma linha otimizada --- como
neste trabalho --- ou um MPC, e o rótulo é o comando limpo que ele emitiria. O
ganho estrutural é o que a subseção anterior antecipou: o especialista precisa de
pose global, a rede treinada não --- ela recebe apenas a varredura e a
proprioceptiva, e produz o comando. A competência que dependia da localização é
\textbf{destilada} numa política que não depende dela.

\subsection{Redes compactas sobre LiDAR bidimensional}

A arquitetura adotada na etapa E2 deste trabalho é a do TinyLidarNet
\cite{Zarrar2024TinyLidarNet}, uma rede convolucional unidimensional que consome
a varredura LiDAR bruta e produz esterço e velocidade. O trabalho original
demonstra que uma rede dessa classe, treinada com um conjunto pequeno de
demonstrações, obteve o terceiro lugar em um Grande Prêmio F1TENTH real, e
enfatiza a viabilidade de inferência em microcontroladores --- que é a
propriedade que interessa a qualquer trabalho com ambição de embarcar a política.

A instância usada aqui é a variante média, com aproximadamente 221 mil
parâmetros, consumindo os 1080 feixes brutos em dois quadros empilhados mais três
grandezas proprioceptivas. Uma decisão de implementação merece registro por sua
relação com a literatura de transferência: o pré-processamento --- tratamento de
leituras não finitas, saturação e normalização --- vive \textbf{dentro} do grafo
da rede, e não numa rotina externa. A consequência é que o artefato exportado
carrega o contrato de observação junto consigo, e o nó que o executa não precisa
manter uma segunda implementação do pré-processamento. É uma resposta direta a
uma classe de defeito que este projeto mediu na própria pele, documentada no
Capítulo~\ref{cap:metodologia}: duas cópias da mesma transformação divergem.

Vale contrastar com \cite{Ivanov2020Verifying}, que treina controladores neurais
sobre 21 a 61 raios com o objetivo explícito de manter o sistema dentro do
alcance da verificação formal. As duas escolhas são coerentes com seus
objetivos e mutuamente exclusivas: reduzir a entrada compra verificabilidade e
custa resolução espacial; consumir a varredura inteira compra desempenho e
custa verificabilidade.

\subsection{O desvio de distribuição e as duas famílias de resposta}

A clonagem de comportamento tem um defeito estrutural conhecido, e é dele que
depende quase todo o desenho da etapa E2. A rede é treinada sobre a distribuição
de estados \textbf{induzida pelo especialista}, mas é executada sobre a
distribuição induzida por ela mesma. Um erro pequeno desloca o veículo para um
estado ligeiramente fora da distribuição de treino, onde o erro seguinte é maior,
e assim por diante. Ross, Gordon e Bagnell \cite{Ross2011DAgger} formalizaram
esse acúmulo mostrando que, no pior caso, o custo cresce quadraticamente com o
horizonte de execução, contra crescimento linear se o problema fosse de
aprendizado supervisionado convencional.

A literatura oferece duas famílias de resposta.

A primeira é \textbf{interativa}: o DAgger \cite{Ross2011DAgger} executa a
política aprendida, coleta os estados que ela de fato visita, e consulta o
especialista para rotular esses estados, agregando o resultado ao conjunto de
treino. A garantia é forte, mas o preço é a exigência de manter o especialista
disponível \textbf{durante} o treino, em cada iteração. Quando o especialista é
humano, isso é caro; quando é um controlador que exige pose global, é uma
dependência que se quer justamente eliminar.

A segunda é \textbf{por injeção de ruído na coleta}: o DART
\cite{Laskey2017DART} perturba a ação \textbf{executada} pelo especialista com
ruído, mas rotula o estado visitado com a ação \textbf{limpa} que o especialista
teria emitido. O veículo passa a visitar estados fora da trajetória ideal ---
que é a distribuição em que a política aprendida vai operar --- e aprende a
retornar dela. A coleta permanece em uma única passada, sem especialista no laço
de treino. É essa a formulação adotada aqui.

Vale registrar, ao lado, o diagnóstico de Codevilla e colaboradores
\cite{Codevilla2019Limitations} sobre os limites da clonagem de comportamento em
condução autônoma: viés do conjunto de dados, confusão causal e alta variância
entre execuções de treino com a mesma configuração. As três aparecem neste
trabalho de forma reconhecível, e o Capítulo~\ref{cap:metodologia} descreve a
instrumentação construída para não confundir nenhuma delas com um resultado.

\subsection{O que a literatura não fixa: a magnitude do ruído}

O DART \cite{Laskey2017DART} formula a escolha da magnitude do ruído como
otimização de um substituto da perda esperada, mas na prática dos trabalhos
aplicados o parâmetro aparece como um escalar único, escolhido por tentativa e
mantido fixo ao longo de toda a coleta. Este trabalho mediu que esse escalar
único é errado por construção quando a coleta atravessa pistas de dificuldade
diferente.

Com o mesmo desvio-padrão de 0,02~rad aplicado às três pistas de coleta, o
especialista percorreu 43~km em Catalunya \textbf{sem uma única colisão},
enquanto colidia a cada 7~km em Spielberg e em Silverstone. O efeito é composto e
perverso: a pista em que o veículo mais colide durante a coleta é a que produz
menos amostras e mais episódios truncados, e é justamente a pista em que a rede
depois falha. Não se trata de incompetência do especialista --- avaliado sem
ruído, ele completa todas as voltas em Silverstone.

A resposta adotada foi calibrar o ruído \textbf{por pista}, contra um alvo de
colisões por quilômetro fixado antes de qualquer medição, percorrendo a grade em
ordem decrescente e escolhendo o maior nível cujo risco medido cabe no alvo ---
porque menos ruído é menos cobertura fora da trajetória, que é o próprio ponto do
DART. A calibração resultante separa Silverstone das demais com metade do ruído.
A métrica é colisões por quilômetro, e não fração de episódios abortados, porque
nesta escala de episódio qualquer risco por passo produz ``metade abortou'', o
que diz mais sobre o comprimento do episódio do que sobre o ruído. A
discussão da incerteza de Poisson associada a poucos eventos por medição --- e do
que torna a escolha do nível robusta apesar dela --- está na
Seção~\ref{sec:val-poisson}.

Esta é uma das cinco diferenças que a Seção~\ref{sec:arte-lacuna} vai
consolidar: \textbf{a calibração por pista, contra um alvo pré-declarado, é uma
operacionalização do DART que a literatura consultada não descreve}.

\subsection{Generalização entre pistas}

Medir uma política de corrida apenas no circuito em que ela foi treinada mede
memorização tanto quanto competência. A literatura da plataforma reconhece o
problema: Bosello, Tse e Pau \cite{Bosello2022CrossTrack} tratam explicitamente
de treinar num circuito e correr em outro com política baseada em LiDAR, e
Brunnbauer e colaboradores \cite{Brunnbauer2022Latent} usam aprendizado por
reforço baseado em modelo com imaginação latente para obter transferência sem
ajuste adicional.

Este trabalho avalia a rede de imitação em circuitos distintos do circuito de
avaliação dos critérios, com um cuidado metodológico que vale explicitar: o
critério é \textbf{relativo ao controle}. Para cada circuito calcula-se a razão
entre a taxa de conclusão da rede e a do próprio especialista, saturada em 1,0, e
excluem-se os circuitos em que o especialista não completa volta alguma. A
justificativa é medida e está no Capítulo~\ref{cap:metodologia}: um limiar
absoluto seria uma trave impossível, já que o especialista completa 0\% das
voltas em Monza e 17\% em Oschersleben; e sem a saturação em 1,0 uma média
poderia ser aprovada com um circuito zerado, compensado por outro em que a rede
supera o especialista.

Nos três circuitos que entram no critério --- Catalunya, Silverstone e
Oschersleben --- a razão medida foi 1,00, tanto na média quanto no pior caso, com
15 de 15 voltas em cada um. Esse número agregado diz menos do que parece, e o que
ele esconde favorece o resultado. Dois dos três circuitos, Catalunya e
Silverstone, integravam o conjunto de coleta: 1,00 ali é empate com o
especialista em pista que ambos percorreram, e não mede generalização. Os dois
circuitos que \textbf{nunca} entraram na coleta são justamente onde a rede mais
se afasta do professor, e o critério não converte em nota a margem de nenhum dos
dois: em Oschersleben a rede completa 15 de 15 voltas contra 1 de 6 do
especialista --- uma razão de aproximadamente \textbf{seis vezes}, que a
saturação em 1,0 corta ---, e em Monza completa 12 de 13 contra 0 de 5, com o
circuito excluído do cálculo porque a razão dividiria por zero.

A afirmação que este trabalho sustenta diante da literatura de imitação é,
portanto, mais forte do que o valor agregado do portão sugere: uma rede destilada
de um especialista geométrico completa \textbf{100\% das voltas em Oschersleben e
92\% em Monza}, dois circuitos que nunca entraram na coleta de demonstrações e
nos quais o próprio especialista que a treinou completa 17\% e 0\%. A destilação
não apenas transfere a competência do especialista para uma política sem pose
global: fora do circuito de coleta ela a \textbf{supera}, num regime em que a
política ensinante falha. O aparato que produz esse número, a sobreposição entre
o conjunto de generalização e as ações de correção do laço automático, e a razão
pela qual o critério de aceite não converte em nota nenhuma das duas margens
estão no Capítulo~\ref{cap:metodologia}.

\section{Aprendizado por reforço para corrida}\label{sec:arte-rl}

\subsection{Agentes treinados do zero e o custo da exploração}

A formulação mais direta do problema por aprendizado por reforço trata a corrida
como um processo de decisão de Markov cuja recompensa premia progresso e pune
colisão, e treina uma política do zero. Os resultados mais impressionantes dessa
linha vêm de domínios com orçamento computacional muito superior ao acadêmico
típico: Wurman e colaboradores \cite{Wurman2022GranTurismo} treinaram um agente
capaz de superar campeões humanos num simulador de corrida comercial, e Kaufmann
e colaboradores \cite{Kaufmann2023DroneRacing} obtiveram desempenho de nível de
campeão em corrida de drones com política treinada em simulação e executada em
voo real.

Na plataforma F1TENTH, a experiência reportada é mais sóbria. Evans, Jordaan e
Engelbrecht \cite{Evans2023Comparing} comparam arquiteturas de aprendizado por
reforço profundo para corrida em escala e documentam a dificuldade de políticas
fim a fim treinadas do zero em obter, simultaneamente, tempo de volta competitivo
e confiabilidade; em trabalho subsequente \cite{Evans2023Trajectory} os mesmos
autores propõem exatamente uma correção desse problema, incorporando uma
trajetória de referência ao sinal de recompensa. O padrão que emerge é
consistente e é a justificativa central da arquitetura adotada aqui: o agente
tabula rasa gasta a maior parte do seu orçamento de amostras aprendendo a não
bater, que é competência que já existe pronta em um controlador clássico.

O FastRLAP \cite{Stachowicz2023FastRLAP} ataca o mesmo problema por outra via, e
é o contraponto mais forte a este trabalho na literatura consultada. O sistema
aprende a dirigir agressivamente em menos de 20 minutos de treino \textbf{online
no mundo real}, sem simulador e sem demonstrações de alto desempenho. Ele
consegue isso combinando três ingredientes: inicialização das representações a
partir de um grande conjunto de dados prévio de navegação, uma única demonstração
de baixa velocidade, e um sistema de prática autônoma com máquina de estados
finitos que seleciona objetivos e se recupera de colisões sem intervenção humana.
As limitações declaradas pelos próprios autores são igualmente relevantes: o
sistema exige um estimador de estado que forneça o vetor para o próximo ponto de
controle, e não trata explicitamente de segurança.

\subsection{Aprendizado por reforço residual}

A formulação residual parte de uma observação simples: se já existe um
controlador que resolve o grosso da tarefa, o aprendizado não precisa
reaprendê-lo. Silver e colaboradores \cite{Silver2018ResidualPolicy} e Johannink
e colaboradores \cite{Johannink2019Residual} propõem, de forma independente e
quase simultânea, decompor a política final em uma base fixa e um resíduo
aprendido, cuja ação se soma à da base:
\[
    a_{\text{final}} = \pi_{\text{base}}(o) + \pi_{\text{res}}(o).
\]

As vantagens são três e todas se aplicam aqui. A exploração parte de uma região
viável do espaço de estados, e não de comportamento aleatório --- desde o
primeiro passo o veículo dá voltas, e a coleta produz transições informativas. A
base absorve a parte da tarefa que a engenharia clássica resolve bem, deixando ao
aprendizado a parte que ela resolve mal --- tipicamente os efeitos difíceis de
modelar. E a decomposição é uma forma de regularização: o resíduo é uma correção,
não uma política completa, e limitar a sua amplitude limita o quanto a política
final pode se afastar do comportamento conhecido da base.

A literatura de controle residual em robótica \cite{Johannink2019Residual}
tipicamente ancora o resíduo num controlador projetado à mão que \textbf{depende
de estado estimado} --- realimentação de posição, controle de impedância,
seguimento de trajetória. É aqui que este trabalho diverge.

\subsection{A escolha deste trabalho: o resíduo sobre a rede de imitação}

Na etapa E3, a ação do agente se soma à saída da \textbf{rede de imitação} da
etapa E2, e não à do especialista de \textit{pure pursuit} da etapa E1:
\[
    a_{\text{final}} = \operatorname{clip}\left(a_{\text{IL}} + k \cdot a_{\text{RL}},\ -1,\ 1\right).
\]

A alternativa --- ancorar o resíduo no \textit{pure pursuit} --- é tentadora e
foi rejeitada deliberadamente. Ela daria uma base mais rápida, já que o
especialista completa a volta em 59,000~s contra 60,050~s da rede que o imita.
Dois argumentos a derrubam.

O primeiro é de \textbf{independência de localização}, e é o argumento decisivo. O
\textit{pure pursuit} exige pose global, portanto exige localização. Se a base
for ele, a política final herda essa dependência inteira: mapa, filtro de
localização, orçamento computacional e modos de falha associados
\cite{fox2003adapting}. A rede de imitação, ao contrário, enxerga apenas LiDAR e
proprioceptiva. Ancorando o resíduo nela, a política final \textbf{não depende de
pose global} e continua usando só o que o sensor entrega --- que é exatamente a propriedade que
justifica a etapa E2 existir, e que seria desperdiçada ao voltar a depender do
especialista na etapa seguinte. A decisão subordina o desempenho da linha de base
à independência de localização, e o faz explicitamente.

O segundo é de \textbf{falha graciosa}. Se o resíduo não aprender nada útil e a
sua saída convergir para zero, a política composta degrada exatamente para a rede
da etapa E2 --- que já passou os seus próprios critérios de aceite. O pior caso
da etapa E3 é o resultado da etapa E2, e não um comportamento indefinido. A
saturação em $[-1, 1]$ mantém a composição dentro do envelope do atuador, e os
ganhos $k$ funcionam como \textbf{limite de autoridade} do resíduo: com
$k_{\text{steer}} = 0{,}15$, a correção de esterço não pode dominar o comando da
base. Reduzir esse ganho pela metade foi, empiricamente, a intervenção que mais
comprou confiabilidade no desenvolvimento da etapa.

O algoritmo é o Soft Actor-Critic \cite{Haarnoja2018SAC}, escolhido por ser fora
de política --- portanto capaz de reaproveitar transições --- e por regularizar a
exploração por entropia, o que é adequado a um espaço de ação contínuo e de baixa
dimensão como o de uma correção. A implementação é a do Stable-Baselines3
\cite{Raffin2021SB3}.

O resultado é o produto final do trabalho: tempo de volta de \textbf{52,525~s},
com intervalo de confiança de 95\% de $[52{,}51;\ 52{,}72]$, o que representa
cerca de \textbf{11\% mais rápido que o especialista} que originou os rótulos, e
7,5~s mais rápido que a rede de imitação que o resíduo corrige. A taxa de
conclusão de voltas medida foi 0,9891, com 91 de 92 voltas completadas, e a
degradação sob randomização do coeficiente de atrito, medida de forma pareada,
foi de 0,0000 em taxa de conclusão e 0,0019 em tempo de volta. A ressalva
estatística que acompanha os três resultados aprovados --- o limiar cai dentro do
intervalo de confiança, por decisão pré-registrada e com motivo escrito antes da
medição --- está no Capítulo~\ref{cap:metodologia} e vale integralmente para este
número.

\subsection{Um custo da formulação residual que a literatura não reporta}

Há um achado deste trabalho que qualifica a formulação residual e que não foi
encontrado na literatura consultada sobre políticas residuais. A política
composta é \textbf{sensível a perturbação numérica irrelevante} num grau que a
política de imitação isolada não é.

Perturbando o vetor de ganhos por valores de $\varepsilon$ que varrem sete ordens
de grandeza e reavaliando o mesmo ponto de controle com as mesmas sementes, a
taxa de conclusão da política composta variou com amplitude de 0,0381, enquanto a
da rede de imitação variou \textbf{zero}. Uma perturbação de $10^{-12}$ desloca a
taxa de conclusão da política composta em quase quatro pontos percentuais. A
perturbação não desloca o resultado numa direção: ela reembaralha \textbf{qual}
episódio termina em colisão.

A explicação é coerente com o que se espera de uma política mais rápida: o
resíduo dirige mais perto do limite, e ali uma diferença mínima decide quem bate.
A consequência para a leitura de qualquer resultado dessa família é que, quando a
amplitude sob perturbação irrelevante é comparável à distância até o limiar, o
número reproduzível é a \textbf{distribuição}, e não a realização --- quem
reimplementar a composição com outra ordem de operações ou outro tipo de ponto
flutuante medirá outro valor dessa faixa. O tratamento completo está no
Capítulo~\ref{cap:metodologia}; o que interessa ao estado da arte é que trabalhos
que reportam um valor pontual de taxa de conclusão para políticas residuais de
alta velocidade podem estar reportando uma realização, e não uma estimativa.

\section{Randomização de domínio}\label{sec:arte-rand}

Treinar e avaliar uma política sob uma única configuração do simulador mede o
quanto ela se ajusta àquela configuração, e não o quanto ela é robusta. A
literatura de transferência entre simulação e realidade consolidou a
\textbf{randomização de domínio} como resposta a isso: Tobin e colaboradores
\cite{Tobin2017DomainRand} randomizam a aparência visual do ambiente simulado,
de modo que o mundo real seja apenas mais uma amostra da distribuição de treino,
e Peng e colaboradores \cite{Peng2018DynamicsRand} aplicam a mesma lógica aos
parâmetros de \textbf{dinâmica}, treinando sob uma distribuição de massas,
atritos e atrasos. Este trabalho adota a ideia pelo seu efeito de robustez, e
não para atravessar qualquer fronteira: o coeficiente de atrito é sorteado a
cada episódio, a largada é sorteada sobre a linha de corrida e o sensor recebe
ruído, nas etapas E2 e E3 e na avaliação de todos os pilotos, e a robustez a
essa variação é um critério de aceite explícito, medido de forma pareada. A
randomização cobre apenas essas três fontes de variação, e a Seção~\ref{sec:arte-lacuna}
registra o que fica de fora.

\section{O formato \textit{time trial} e a comparação com pilotos humanos}\label{sec:arte-timetrial}

A comparação de tempo de volta feita neste trabalho imita o formato \textit{time
trial} da competição F1TENTH. Nessa modalidade, cada veículo percorre a pista
individualmente, em duas baterias de 5 minutos, buscando completar o maior número
possível de voltas consecutivas sem interrupções e registrar a volta mais rápida
do período alocado.

A pontuação é deliberadamente construída para equilibrar velocidade e
consistência: soma os pontos do \emph{ranking} dos tempos de volta mais rápidos
aos pontos do \emph{ranking} do número de voltas consecutivas ininterruptas
\cite{F1TENTHRules2023}. Isso força quem compete a balancear trajetórias
agressivas, que dão velocidade máxima, com controle que se mantenha consistente
ao longo de várias voltas e que minimize a probabilidade de falha. É esse mesmo
par de grandezas, volta rápida e voltas completas sem colisão, que as métricas de
conclusão e de tempo de volta deste trabalho medem, para todos os pilotos sob o
mesmo avaliador.

\section{Pilotos humanos e controladores autônomos em corrida}\label{sec:arte-humanos}

Comparar um controlador com um piloto humano é uma prática antiga da área, mas
raramente é o objeto central do estudo: em geral o humano aparece como régua de
um resultado cujo foco é o algoritmo. Esta seção reúne o que esses trabalhos
reportam e delimita o que deles se aproveita.

\textbf{Desempenho humano como referência.} Os resultados mais citados usam o
humano como teto a superar. Wurman e colaboradores \cite{Wurman2022GranTurismo}
medem o agente contra campeões humanos em um simulador comercial, e Kaufmann e
colaboradores \cite{Kaufmann2023DroneRacing} fazem o mesmo em corrida de drones.
Fuchs e colaboradores \cite{Fuchs2021SuperHuman} comparam o tempo de volta de uma
política treinada por aprendizado por reforço com os de uma base de mais de
50\,000 jogadores do mesmo simulador, e reportam que ela supera o mais rápido da
base. Nesses trabalhos o humano é um especialista com muita prática e a comparação
central é de tempo de volta. Igualar a informação do humano à do agente e comparar a
suavidade do comando não são o foco deles; são esses dois eixos que o presente
trabalho acrescenta, em escala bem menor.

\textbf{Imitação de pilotos humanos.} A clonagem de comportamento, vista na
Seção~\ref{sec:arte-il}, tem na condução humana a sua origem
\cite{Pomerleau1989ALVINN, Bojarski2016EndToEnd}. Em corrida, Löckel, Peters e
van Vliet \cite{Lockel2020ProMoD} propõem um arcabouço probabilístico para imitar
o comportamento de pilotos de corrida em simulação, com trajetória-alvo e rede
que escolhe as ações. Este trabalho não imita humanos: o especialista que
gera as demonstrações do E2 é geométrico, e a imitação de demonstrações humanas
fica como trabalho futuro.

\textbf{Aprendizado do piloto ao longo das tentativas.} Um humano melhora com a
prática, e o tempo de volta numa sessão curta mistura habilidade com
familiaridade. Por isso a medida de uma pessoa depende do treino que ela teve
antes de ser medida, e um protocolo de comparação precisa fixar esse treino e
registrar a evolução do tempo de volta, em vez de apresentar um único número.
Esta é uma consequência do modo como o aprendizado humano é medido em geral, e
não uma afirmação extraída dos trabalhos citados acima, que não a examinam.

\textbf{Informação disponível ao piloto.} O desempenho de quem pilota depende do
que ele vê. Uma pessoa com o mapa da pista pode antecipar curvas e escolher
velocidade com antecedência; uma política que recebe apenas a varredura do LiDAR
no passo atual decide só pelo que está à frente. Comparar um humano com a vista
da pista inteira contra essa política mistura habilidade com informação. Por isso
o comparativo mede o humano também com a vista restrita ao LiDAR, em que a
informação é equivalente à da política, e declara a assimetria da primeira vista
(Seção~\ref{sec:val-limitacoes-comparativo}).

\section{Posicionamento: a lacuna e os seus limites}\label{sec:arte-lacuna}

\subsection{O que este trabalho faz que os trabalhos citados não fazem}

Cinco diferenças sustentam a contribuição, e nenhuma delas é ``um algoritmo
novo''. Todas são decisões de composição e de método. A primeira diz respeito ao
comparativo, que é o objeto do trabalho; as outras quatro, à cascata de pilotos
aprendidos.

\textbf{Primeira: os três tipos de piloto são medidos sob o mesmo avaliador e com
a informação igualada.} Os trabalhos que comparam políticas com humanos
(Seção~\ref{sec:arte-humanos}) comparam sobretudo tempo de volta, com o humano na
vista que o jogo lhe oferece. Aqui controladores clássicos, políticas aprendidas
e pilotos humanos correm no mesmo simulador, com os mesmos limites de ação, a
mesma randomização e o mesmo avaliador, e o humano é medido também na vista só
LiDAR, em que recebe informação equivalente à da política. Além do tempo de volta,
a comparação cobre conclusão, colisões, suavidade do esterço, desvio da linha de
corrida e perfil de velocidade.

\textbf{Segunda: o resíduo é ancorado numa política que só enxerga LiDAR.} A
literatura de aprendizado por reforço residual ancora o resíduo em controladores
que dependem de estado estimado \cite{Silver2018ResidualPolicy,
Johannink2019Residual}. O BayesRace \cite{Jain2020Bayesrace}, que é o vizinho
conceitual mais próximo, coloca o resíduo no espaço do \textbf{modelo}, corrigido
por processos gaussianos e consumido por um MPC que exige pose e solucionador.
Aqui a base é uma rede de imitação sem pose global, o resíduo vive no espaço da
\textbf{ação}, e a composição inteira continua dependendo só do sensor. A escolha custa
desempenho de linha de base --- 60,050~s contra 59,000~s do especialista --- e
compra independência de localização. É um compromisso explícito, e não uma
preferência estética: a política passa a receber a mesma informação que um
piloto humano com a vista restrita ao LiDAR.

\textbf{Terceira: falha graciosa por construção.} Com a saturação e o limite de
autoridade do resíduo, o pior caso da política composta é o comportamento da
política de imitação, que já satisfaz critérios próprios. Essa propriedade é
mencionada qualitativamente na literatura residual; aqui ela é uma consequência
verificável da forma da composição, e a política de imitação que serve de piso
tem os seus próprios números medidos.

\textbf{Quarta: o ruído de coleta é calibrado por pista contra um alvo
pré-declarado.} O DART \cite{Laskey2017DART} fixa o arcabouço; a prática aplicada
usa um escalar único. Este trabalho mediu que um escalar único produz risco de
colisão que varia por um fator de ordem dez entre circuitos, e que o circuito mais
arriscado é o que produz menos dado e é o que a rede depois falha. A calibração
resultante é reportada com o alvo, a grade percorrida e o risco medido em cada
nível.

\textbf{Quinta: os critérios de aceite são congelados antes da medição e viajam
com o experimento.} A prática corrente na área reporta métrica agregada mais
descrição de treino \cite{Betz2022Survey}. Este trabalho trata o momento em que o
critério foi fixado como parte do resultado: um resultado obtido com critérios
que não coincidem com os congelados no nascimento do experimento é declarado
inválido, e não aprovado nem reprovado. O aparato completo é o assunto do
Capítulo~\ref{cap:metodologia}, e é ele que dá sentido à sexta observação, que
não é uma contribuição mas uma ressalva: a medição da sensibilidade numérica da
política composta, reportada como limitação em vez de omitida.

A Tabela~\ref{tab:arte_posicionamento} resume o posicionamento em relação aos
trabalhos mais próximos.

\begin{table}[htb]
    \centering
    \small
    \begin{tabular}{|p{3.1cm}|p{2.3cm}|p{2.6cm}|p{2.2cm}|p{2.4cm}|}
        \hline
        \textbf{Trabalho} & \textbf{Paradigma} & \textbf{Política final exige pose global?} &
        \textbf{Validação} & \textbf{Critério congelado antes?} \\
        \hline
        \cite{Wang2020LMPC} & LMPC & Sim & Simulação & Não declarado \\
        \hline
        \cite{Jain2020Bayesrace} & MPC + resíduo no modelo & Sim & Simulação (duas escalas) & Não declarado \\
        \hline
        \cite{Stachowicz2023FastRLAP} & RL online no mundo real & Sim (vetor ao próximo alvo) & Hardware & Não declarado \\
        \hline
        \cite{Zarrar2024TinyLidarNet} & Imitação fim a fim & Não & Hardware (competição) & Não declarado \\
        \hline
        \cite{Ivanov2020Verifying} & RL + verificação formal & Não & Hardware + prova & Propriedade formal \\
        \hline
        \cite{Evans2023Trajectory} & RL com recompensa guiada por trajetória & Trajetória de referência exigida no treino & Simulação & Não declarado \\
        \hline
        Este trabalho & Imitação + resíduo de RL & \textbf{Não} & \textbf{Simulação} & \textbf{Sim} \\
        \hline
    \end{tabular}
    \caption{Posicionamento em relação aos trabalhos mais próximos. A coluna de
    critério congelado registra se o texto declara que os limiares de aceite
    foram fixados antes da medição.}
    \label{tab:arte_posicionamento}
\end{table}

A tabela não inclui os trabalhos da Seção~\ref{sec:arte-humanos}, porque as suas
colunas (exigência de pose global, validação em hardware, critério congelado)
foram desenhadas para políticas aprendidas sem pose, e os trabalhos de comparação
com humanos citados ali não reportam esses campos de forma que permita preencher
a linha sem suposição.

\subsection{O que os trabalhos citados fazem e este não faz}

Um estado da arte que só elogia a literatura não posiciona nada. As lacunas deste
trabalho em relação aos citados são as seguintes, e são declaradas sem atenuação.

\begin{itemize}
    \item \textbf{Nenhum experimento em hardware.} O FastRLAP
          \cite{Stachowicz2023FastRLAP}, o TinyLidarNet
          \cite{Zarrar2024TinyLidarNet}, o trabalho de verificação
          \cite{Ivanov2020Verifying}, a transferência de
          \cite{Chisari2021LearningSim} e \cite{Brunnbauer2022Latent} e a
          própria plataforma \cite{OKelly2020F1TENTH} medem no hardware. Este
          trabalho mede em simulação, por escolha de escopo.
          \textbf{Toda} comparação numérica deste texto com resultados de campo
          é, portanto, assimétrica.
    \item \textbf{Nenhuma corrida direta nem ultrapassagem.} O regulamento da
          categoria \cite{F1TENTHRules2023} prevê corrida com dois veículos, e a
          literatura trata do assunto tanto do lado do planejamento quanto do
          lado do teste adversarial \cite{Bak2021Stress, Betz2022Survey}. Este
          trabalho ataca exclusivamente o cenário de veículo único.
    \item \textbf{Nenhuma garantia formal.} A abordagem de
          \cite{Ivanov2020Verifying} produz uma prova sobre a política; aqui há
          um critério estatístico sobre uma amostra finita, e a própria escolha
          de arquitetura --- rede convolucional sobre 1080 feixes --- coloca a
          política fora do alcance das ferramentas de verificação descritas.
    \item \textbf{Nenhuma adaptação de modelo em linha.} LMPC
          \cite{Rosolia2020LMPC, Wang2020LMPC}, MPC com processos gaussianos
          \cite{Kabzan2019LearningMPC} e BayesRace \cite{Jain2020Bayesrace}
          melhoram a cada volta executada. A política aqui é fixa após o treino.
    \item \textbf{Nenhuma percepção por câmera nem fusão sensorial.} A política
          consome exclusivamente LiDAR bidimensional e proprioceptiva. Isso é uma
          escolha coerente com a independência de localização, mas exclui o trabalho da
          discussão de percepção em alta velocidade que a revisão da área
          identifica como desafio central \cite{Betz2022Survey}.
    \item \textbf{Um simulador, um modelo de veículo, um circuito de portão.} A
          fidelidade do modelo de faixa única e a representatividade do circuito
          de avaliação são premissas herdadas, não resultados verificados. A
          Seção~\ref{sec:val-colisao} lista as limitações do detector de
          colisão do simulador que precisaram ser corrigidas antes que qualquer
          número fosse reportável --- o que é, por si, evidência de que a premissa
          de fidelidade merece a desconfiança que este parágrafo lhe dá.
\end{itemize}

\subsection{Síntese}

A literatura de corrida autônoma em escala oferece, de um lado, controle clássico
que alcança desempenho próximo ao ótimo ao preço de um modelo, um estimador de
estado e um solucionador \cite{Heilmeier2020MinCurv, Wang2020LMPC,
Jain2020Bayesrace}; e, de outro, políticas aprendidas fim a fim que dispensam
tudo isso ao preço de amostras, de garantias e de previsibilidade
\cite{Evans2023Comparing, Stachowicz2023FastRLAP}. O aprendizado por imitação
ocupa o meio-termo, destilando a competência do primeiro grupo numa política com
as dependências do segundo, e paga por isso com o desvio de distribuição que o
DART \cite{Laskey2017DART} e o DAgger \cite{Ross2011DAgger} atacam por caminhos
distintos.

A cascata de pilotos aprendidos deste trabalho percorre esse arranjo em três passos, e cada
passo preserva a propriedade conquistada no anterior: o especialista geométrico
fornece competência; a imitação a transfere para uma política sem pose global; e
o resíduo de aprendizado por reforço recupera --- e supera --- o desempenho
perdido na destilação, sem devolver a dependência de localização e sem que o pior
caso seja pior que a política que ele corrige. O resultado, 52,525~s contra
59,000~s do especialista, é a medida dessa recuperação.

O que falta a essa literatura, e o que este trabalho se propõe a medir, é a
posição desses pilotos entre si e em relação a uma referência humana, no mesmo
simulador e sob os mesmos critérios.

Os demais pilotos do comparativo ocupam as pontas desse arranjo. O Disparity
Extender e o seguidor de gap, descritos entre os métodos reativos
(Seção~\ref{sec:arte-classico}), são os reativos sem mapa: decidem só pela
varredura atual, como a rede, mas sem aprender nada. O
piloto humano é a referência de desempenho e de modo de decidir
(Seção~\ref{sec:arte-humanos}), medida sob o formato \textit{time trial} da
Seção~\ref{sec:arte-timetrial}: volta rápida e voltas completas sem colisão.
